\documentclass[11pt]{article}
% ======================================================================
%  examstyle.tex — EPFL-style exam layout for the weekly Time Series exams
%  Shared by every Exam_Week_NN(.tex / _Solutions.tex). Compiles with pdflatex.
% ======================================================================
\usepackage[utf8]{inputenc}
\usepackage[T1]{fontenc}
\usepackage[english]{babel}
\usepackage[a4paper,margin=2.4cm]{geometry}
\usepackage{amsmath,amssymb,amsthm,mathtools}
\usepackage{bm}
\usepackage{enumitem}
\usepackage{array}

\setlength{\parindent}{0pt}
\setlength{\parskip}{0.5em}
\emergencystretch=3em
\allowdisplaybreaks[1]

% ---------- math macros (same names as the rest of the project) ----------
\newcommand{\E}{\mathbb{E}}
\DeclareMathOperator{\Var}{Var}
\DeclareMathOperator{\Cov}{Cov}
\DeclareMathOperator{\Corr}{Corr}
\DeclareMathOperator*{\argmin}{arg\,min}
\DeclareMathOperator{\MSE}{MSE}
\DeclareMathOperator{\trace}{tr}
\newcommand{\Reals}{\mathbb{R}}
\newcommand{\Ints}{\mathbb{Z}}
\newcommand{\Comp}{\mathbb{C}}
\newcommand{\Nat}{\mathbb{N}}
\newcommand{\Prob}{\mathbb{P}}
\newcommand{\Tr}{^{\mathsf{T}}}
\newcommand{\conj}[1]{\overline{#1}}
\newcommand{\abs}[1]{\left\lvert #1 \right\rvert}
\newcommand{\norm}[1]{\left\lVert #1 \right\rVert}
\newcommand{\ind}{\mathbf{1}}
\newcommand{\eps}{\varepsilon}
\newcommand{\diff}{\nabla}
\newcommand{\acvs}{\gamma}
\newcommand{\acf}{\rho}
\newcommand{\pacf}{\alpha}
\newcommand{\sdf}{\mathcal{S}}
\newcommand{\Bsh}{B}
\newcommand{\iid}{\overset{\text{iid}}{\sim}}

\setlist[enumerate]{leftmargin=2em,topsep=2pt,itemsep=4pt}
\setlist[itemize]{leftmargin=1.6em,topsep=2pt,itemsep=2pt}

\newcounter{exo}

% ---------- exam cover header :  \examheader{exam title}{duration} ----------
\newcommand{\examheader}[2]{%
  \thispagestyle{empty}
  \begin{center}
    {\Large\bfseries Time Series}\\[3pt]
    Spring Semester 2026, EPFL\\
    \'Ecole Polytechnique F\'ed\'erale de Lausanne\\
    Prof.\ Dr.\ Sofia C.\ Olhede\\[7pt]
    \hrule height 0.8pt
    \vspace{9pt}
    {\large\bfseries Practice Exam --- #1}\\[4pt]
    {\normalsize Duration: #2 \quad$\cdot$\quad No calculator \quad$\cdot$\quad Answer on paper}\\
    \vspace{8pt}
    \hrule height 0.8pt
  \end{center}
  \vspace{14pt}
  \begin{tabular}{@{}p{8.2cm}@{}p{8cm}@{}}
    Name:\enspace\hrulefill & Surname:\enspace\hrulefill
  \end{tabular}\\[14pt]
  SCIPER (Student Card Number):\enspace\hrulefill
  \vspace{16pt}
}

% ---------- points table (fixed: 4 problems) ----------
\newcommand{\pointstable}{%
  \begin{center}
  \renewcommand{\arraystretch}{1.5}
  \begin{tabular}{|p{5cm}|p{4cm}|}
    \hline
    \textbf{Exercise} & \textbf{Points}\\\hline
    1 & \\\hline
    2 & \\\hline
    3 & \\\hline
    4 & \\\hline
    \textbf{Total} & \\\hline
  \end{tabular}
  \end{center}
}

% ---------- problem heading (exam: one problem per page) ----------
\newcommand{\problem}[1]{%
  \clearpage
  \stepcounter{exo}%
  \begin{center}\textbf{\large Problem \theexo\ (#1 points)}\end{center}
  \vspace{4pt}
}
% ---------- answer space: fill the statement page + one blank answer page ----------
% Put \answerpage at the END of each problem so every exercise gets its own page(s).
\newcommand{\answerpage}{\par\vfill\newpage\null\newpage}

% ---------- solutions document ----------
\newcommand{\solheader}[1]{%
  \begin{center}
    {\Large\bfseries Time Series --- Solutions}\\[3pt]
    {\large Practice Exam --- #1}\\[2pt]
    EPFL, MATH-342 \quad$\cdot$\quad Prof.\ S.\ C.\ Olhede\\[5pt]
    \hrule height 0.8pt
  \end{center}
  \vspace{10pt}
}
\newcommand{\solproblem}[1]{%
  \bigskip
  \stepcounter{exo}%
  {\large\bfseries Problem \theexo\ (#1 points)}\par\nobreak\vspace{2pt}
}
% Rappel de l'énoncé avant la solution (dans les corrigés)
\newcommand{\statementstart}{\par\vspace{1pt}{\bfseries\itshape Statement.}\par\nobreak\begingroup\small\itshape}
\newcommand{\statementend}{\par\endgroup\vspace{5pt}\hrule\vspace{6pt}{\bfseries Solution.}\par\nobreak\vspace{2pt}}

\begin{document}

\examheader{Week 5: Trends, Seasonality \& SARIMA}{60 minutes}
\pointstable
\begin{center}\itshape Justify every answer. Throughout, $\eps_t$ is a mean-zero white noise of variance $\sigma^2$.\end{center}

% ---------------------------------------------------------------
\problem{5}
Consider the signal-plus-noise model $X_t=\mu_t+Y_t$, where $\mu_t$ is a deterministic trend and $\{Y_t\}$ is a zero-mean stationary process with autocovariance sequence $\acvs_\tau$. Recall the difference operator $\diff=I-\Bsh$ and the seasonal difference $\diff_{(s)}=I-\Bsh^{s}$.

\textbf{(a)} Take $\mu_t=a+bt$ (linear trend). Show that $\diff X_t=b+\diff Y_t$, and explain in one sentence why $\diff Y_t$ is stationary while $\{X_t\}$ in general is not. \hfill(1.5 pts)

\textbf{(b)} Now let $\mu_t=a+bt+s_t$ with a \emph{quarterly} season $s_{t+4}=s_t$. Compute $\diff\diff_{(4)}X_t$ \emph{explicitly} as a linear combination of the $Y$'s, and verify that its mean is $0$ and that it is weakly stationary. \hfill(2 pts)

\textbf{(c)} \emph{(A classic trap.)} The operators $\diff_{(s)}=I-\Bsh^{s}$ and $\diff^{s}=(I-\Bsh)^{s}$ are \emph{not} the same. For the deterministic sequence $\mu_t=t^{2}$, compute both $\diff^{2}\mu_t$ and $\diff_{(2)}\mu_t$ and state which kind of structure each operator removes. \hfill(1.5 pts)
\answerpage

% ---------------------------------------------------------------
\problem{5}
Let $\{X_t\}$ be the multiplicative seasonal process $\mathrm{SARMA}(0,1)\times(1,0)_{4}$ (quarterly data), defined by
\[
\bigl(1-\Phi\,\Bsh^{4}\bigr)X_t=\bigl(1-\theta\Bsh\bigr)\eps_t,
\qquad \abs\Phi<1 ,
\]
a seasonal autoregression at lag $4$ driven by a non-seasonal $\mathrm{MA}(1)$ input.

\textbf{(a)} Write the explicit one-step recursion for $X_t$, and by inverting the seasonal AR factor obtain the $\mathrm{MA}(\infty)$ weights $\psi_j$ in $X_t=\sum_{j\ge0}\psi_j\eps_{t-j}$ (give a closed form for $\psi_{4l}$ and $\psi_{4l+1}$, and state that all other $\psi_j=0$). \hfill(2 pts)

\textbf{(b)} Using the white-noise filter rule $\acvs_\tau=\sigma^{2}\sum_{j\ge0}\psi_j\psi_{j+\tau}$, compute $\acvs_0,\acvs_1,\acvs_3,\acvs_4,\acvs_5$ in closed form, and explain why the autocovariance lives in clusters around the seasonal lags $0,4,8,\dots$ \hfill(2 pts)

\textbf{(c)} Evaluate the autocorrelations $\acf_1,\acf_3,\acf_4$ numerically for $\Phi=\tfrac12$, $\theta=\tfrac25$, $\sigma^{2}=1$. \hfill(1 pt)
\answerpage

% ---------------------------------------------------------------
\problem{5}
This problem studies non-stationary integrated models.

\textbf{(a)} \emph{(Random walk.)} Let $X_t=\sum_{j=1}^{t}\eps_j$ with $X_0=0$ (an $\mathrm{ARIMA}(0,1,0)$). Compute $\E[X_t]$, $\Var(X_t)$ and $\Cov(X_t,X_{t+\tau})$ for $\tau\ge0$, deduce that $\{X_t\}$ is \emph{not} stationary, and show $\diff X_t$ is white noise. Show also that $\Corr(X_t,X_{t+\tau})\to1$ as $t\to\infty$ with $\tau$ fixed. \hfill(2 pts)

\textbf{(b)} \emph{(IMA$(1,1)$.)} Let $Y_t=Y_{t-1}+\eps_t-\theta\eps_{t-1}$ with $Y_0=0$. By unrolling the recursion, show that for $t\ge1$
\[
Y_t=\eps_t+(1-\theta)\sum_{j=1}^{t-1}\eps_j-\theta\eps_0 ,
\]
and hence compute $\Var(Y_t)$. For $\theta\neq1$, what is the order of growth of $\Var(Y_t)$ in $t$? \hfill(2 pts)

\textbf{(c)} \emph{(Over-differencing.)} What happens to $Y_t$ and to $\Var(Y_t)$ at the boundary $\theta=1$? Show that $\diff Y_t$ is then a \emph{non-invertible} $\mathrm{MA}(1)$, and state the value of its lag-$1$ autocorrelation $\acf_1$. Why is this the warning signature of over-differencing? \hfill(1 pt)
\answerpage

% ---------------------------------------------------------------
\problem{5}
\textbf{(Revision of Week 4: Yule--Walker estimation.)} Let $\{X_t\}$ be a mean-zero causal $\mathrm{AR}(2)$ process $X_t=\phi_1X_{t-1}+\phi_2X_{t-2}+\eps_t$ observed at $X_1,\dots,X_n$, with sample autocovariances
\[
\hat\acvs_0=10,\qquad \hat\acvs_1=5,\qquad \hat\acvs_2=1 .
\]

\textbf{(a)} Derive the Yule--Walker equations $\acvs_k=\phi_1\acvs_{k-1}+\phi_2\acvs_{k-2}$ ($k\ge1$) and the variance identity $\acvs_0=\phi_1\acvs_1+\phi_2\acvs_2+\sigma^{2}$, making clear where \emph{causality} is used. \hfill(2 pts)

\textbf{(b)} Set up the $2\times2$ Yule--Walker system, invert it \emph{by hand}, and report the estimates $\hat\phi_1,\hat\phi_2$ and $\hat\sigma^{2}$. \hfill(2 pts)

\textbf{(c)} Verify that the fitted model is causal (stationary), e.g.\ by locating the roots of $\hat\Phi(z)$ or via the $\mathrm{AR}(2)$ stationarity triangle. \hfill(1 pt)
\answerpage

\end{document}
